\MethodStyle
\PartDivider{PART V}{CATRA Part B: TCP Rate Adaptation}{Use MAC/PHY utilization to regulate per-flow transport demand.}

\begin{frame}{Why MAC Fairness Is Not Enough}
\begin{center}
\begin{tikzpicture}[node distance=1.0cm]
\node[draw=MethodColor,circle,fill=CardBg] (s2) {S2};
\node[draw=MethodColor,circle,fill=CardBg,right=of s2] (s1) {S1};
\node[draw=MethodColor,circle,fill=CardBg,right=of s1] (r) {R};
\draw[-{Latex[length=2mm]},draw=MethodColor] (s2)--node[above,font=\fontsize{9}{10.5}\selectfont]{Flow 2}(s1);
\draw[-{Latex[length=2mm]},draw=MethodColor] (s1)--node[above,font=\fontsize{9}{10.5}\selectfont]{Flows 1 + 2}(r);
\end{tikzpicture}
\end{center}
\vspace{0.25cm}
\begin{itemize}
\item Channel access control can improve the Fair Bandwidth Ratio between stations.
\item It does not necessarily ensure fair throughput for each TCP flow.
\item CATRA therefore controls TCP sending rate to achieve per-flow fairness.
\end{itemize}
\vspace{0.35cm}
\KeyPoint{Part A improves station access; Part B controls TCP sending rate for fair throughput among flows.}
\end{frame}

\begin{frame}{Flow-Level Bandwidth Ratios}
\small
\begin{columns}[T,onlytextwidth]
\column{0.48\textwidth}
\textbf{Fair share}
\[
\boxed{FBR_f=\frac{1}{n_{total}}}
\]
\begin{itemize}
\item Equal share using the same local $n_{total}$ model as Part A.
\item Same-source/destination TCP traffic is treated as one flow.
\end{itemize}
\column{0.48\textwidth}
\textbf{Real share}
\[
\boxed{RBR_f=\frac{T_f^{Active}}{EP}}
\]
\begin{itemize}
\item Airtime attributed to the examined TCP flow.
\item Measured during the Estimation Period.
\end{itemize}
\end{columns}
\vspace{0cm}
\KeyPoint{Algorithm 2 compares $RBR_f/FBR_f$ with $High_{th}$ and $Low_{th}$ to select the TCP rate adjustment.}
\end{frame}

\begin{frame}{Current TCP Sending Demand}
\PaperClaim{The paper estimates how many packets are ready to send from TCP sequence and window state.}
\[
\boxed{win=cwnd+highest\_ack-cur\_seqno}
\]
\begin{tabular}{p{3.2cm}p{7.8cm}}
\toprule
\textbf{Variable} & \textbf{Meaning in the paper}\\
\midrule
$cwnd$ & TCP congestion-window size.\\
$highest\_ack$ & Highest sequence number acknowledged by a received TCP ACK.\\
$cur\_seqno$ & Last sequence number of a transmitted TCP DATA packet.\\
$win$ & Packets currently ready to be transmitted.\\
\bottomrule
\end{tabular}
\end{frame}

\begin{frame}{Per-Flow Packet Transmission Time}
\PaperClaim{MAC observation supplies both active time and the number of transmitted packets for flow $f$.}
\[
\boxed{T_f^{tr}=\frac{T_f^{Active}}{N_f}}
\]
\begin{columns}[T,onlytextwidth]
\column{0.48\textwidth}
\textbf{$T_f^{Active}$}
\begin{itemize}
\item Average time for transmitting packets of TCP flow $f$ during $EP$.
\end{itemize}
\column{0.48\textwidth}
\textbf{$N_f$}
\begin{itemize}
\item Number of packets from TCP flow $f$ transmitted during $EP$.
\end{itemize}
\end{columns}
\end{frame}

\begin{frame}{Time Required to Transmit the Ready Window}
\PaperClaim{At its fair share, flow $f$ should spread its ready packets over a time proportional to the number of competing flows.}
\[
\boxed{T_f=n_{total}\cdot win\cdot T_f^{tr}}
\]
\begin{center}
\begin{tikzpicture}[node distance=0.65cm]
\node[draw=MethodColor,rounded corners,fill=CardBg] (n) {$n_{total}$ competitors};
\node[draw=MethodColor,rounded corners,fill=CardBg,right=of n] (w) {$win$ packets};
\node[draw=MethodColor,rounded corners,fill=CardBg,right=of w] (t) {$T_f^{tr}$ per packet};
\node[draw=ConclusionColor,rounded corners,fill=CardBg,right=of t] (tf) {$T_f$ fair duration};
\draw[-{Latex[length=2mm]}] (n)--(w);
\draw[-{Latex[length=2mm]}] (w)--(t);
\draw[-{Latex[length=2mm]}] (t)--(tf);
\end{tikzpicture}
\end{center}
\vspace{0.35cm}
\SmallNote{Equivalent derivation in the paper: $T_f=(win\,s_f)/(FBR_f B)$.}
\end{frame}

\begin{frame}{Rate-Control Delay}
\PaperClaim{CATRA slows an over-allocated flow by delaying new packet generation.}
\[
\boxed{\Delta_f=\frac{RBR_f}{FBR_f}T_f}
\]
\begin{center}
\resizebox{0.96\textwidth}{!}{%
\begin{tikzpicture}[node distance=0.58cm]
\node[draw=WarningColor,rounded corners,fill=CardBg] (over) {$RBR_f>FBR_f$};
\node[draw=WarningColor,rounded corners,fill=CardBg,right=of over] (delay) {Larger $\Delta_f$};
\node[draw=MethodColor,rounded corners,fill=CardBg,right=of delay] (space) {Space packet generation};
\node[draw=ConclusionColor,rounded corners,fill=CardBg,right=of space] (chance) {Neighbors gain airtime};
\draw[-{Latex[length=2mm]}] (over)--(delay);
\draw[-{Latex[length=2mm]}] (delay)--(space);
\draw[-{Latex[length=2mm]}] (space)--(chance);
\end{tikzpicture}
}
\end{center}
\vspace{0.4cm}
\begin{itemize}
\item The delay is applied before generating new packets after the current $win$ packets complete transmission.
\item Algorithm 2 also decreases $cwnd$ when the flow is over-allocated and increases it when the flow is under-allocated.
\end{itemize}
\end{frame}

\begin{frame}{Worked Example: Flow 1 at S1}
\begin{columns}[T,onlytextwidth]
\column{0.46\textwidth}
\textbf{Initial bandwidth ratio}
\[
n_{total}=3,\qquad FBR_f=\frac13
\]
\[
RBR_f=\frac12>FBR_f
\]
\column{0.50\textwidth}
\textbf{Applied delay}
\[
\Delta_f=\frac{1/2}{1/3}T_f=\frac32T_f
\]
\[
\Delta_f=\frac92\,win\,T_f^{tr}.
\]
\end{columns}
\vspace{0.12cm}
\begin{center}
\[
RBR'_f=\frac{2}{9},\qquad
\frac12\left(RBR_f+RBR'_f\right)
=\frac12\left(\frac12+\frac29\right)=\frac{13}{36}\approx\frac13=FBR_f
\]
\end{center}
\vspace{-0.18cm}
\KeyPoint{The delay moves Flow 1's average bandwidth ratio toward its fair share and gives neighboring flows more transmission opportunities.}
\end{frame}

\begin{frame}{Algorithm 2: TCP Rate Adaptation}
\begin{columns}[T,onlytextwidth]
\column{0.57\textwidth}
\begin{tikzpicture}[
node distance=0.31cm,
box/.style={draw=MethodColor,rounded corners,fill=CardBg,text width=6.3cm,align=left,inner xsep=7pt,inner ysep=5pt,font=\fontsize{9}{10.5}\selectfont},
flow/.style={-{Latex[length=2mm]},draw=MethodColor}
]
\node[box] (ack) {For each newly received TCP ACK, compute $\rho_f=RBR_f/FBR_f$.};
\node[box,below=of ack] (high) {If $\rho_f>High_{th}$: set $\Delta_f=\rho_fT_f$ and $cwnd\leftarrow\max(cwnd-1,1)$.};
\node[box,below=of high] (low) {Else if $\rho_f<Low_{th}$: set $\Delta_f=0$ and $cwnd\leftarrow cwnd+1$.};
\node[box,below=of low] (orig) {Otherwise: call the original TCP mechanism.};
\draw[flow] (ack)--(high);
\draw[flow] (high)--(low);
\draw[flow] (low)--(orig);
\end{tikzpicture}
\column{0.39\textwidth}
\textbf{Three control regions}
\begin{itemize}
\item \textbf{Over-allocated:} delay and decrease demand.
\item \textbf{Under-allocated:} remove delay and increase demand.
\item \textbf{Near target:} preserve native TCP behavior.
\end{itemize}
\vspace{0.2cm}
\SmallNote{Parameters: $High_{th}=1.05$, $Low_{th}=0.7$, $EP=2$ s.}
\end{columns}
\end{frame}
